% Generated draft: requires mathematical and source review.
\documentclass[11pt]{article}
\usepackage[margin=25mm]{geometry}
\usepackage{fontspec}
\setmainfont{Latin Modern Roman}
\usepackage{amsmath,amssymb,mathtools,unicode-math}
\setmathfont{Latin Modern Math}
\usepackage{enumitem,needspace,xurl,hyperref}
\setlength{\parindent}{0pt}
\setlength{\parskip}{5pt}
\setlength{\emergencystretch}{3em}
\newcommand{\N}{\mathbb{N}}
\newcommand{\Z}{\mathbb{Z}}
\newcommand{\Q}{\mathbb{Q}}
\newcommand{\R}{\mathbb{R}}
\newcommand{\C}{\mathbb{C}}
\newcommand{\F}{\mathbb{F}}
\newcommand{\A}{\mathbb{A}}
\newcommand{\PP}{\mathbb{P}}
\newcommand{\E}{\mathbb{E}}
\newcommand{\eps}{\varepsilon}
\newcommand{\dd}{\,\mathrm{d}}
\begin{document}
\section*{UnsolvedMath ID: 800040}
\textbf{Source problem code:} catalogue entry
\subsection{Definitions and mathematical statement}
For a semigroup \(S\), write its multiplication as \(\cdot\); thus \((x\cdot y)\cdot z=x\cdot(y\cdot z)\) for all \(x,y,z\in S\).

The supplied record uses the technical term \emph{identical inclusion} without defining its syntax or satisfaction relation. Relative to the intended source definition, let \(\mathcal L_{\mathrm{inc}}\) denote the collection of identical inclusions and put
\[
\operatorname{Inc}(S)=\{\varphi\in\mathcal L_{\mathrm{inc}}:S\models\varphi\}.
\]
Two semigroups \(S\) and \(T\) are called \emph{inclusively equivalent} if
\[
S\sim_{\mathrm{inc}}T
\quad\Longleftrightarrow\quad
\operatorname{Inc}(S)=\operatorname{Inc}(T).
\]
The source also states that this is equivalent to \(S\) and \(T\) generating the same \emph{inclusive variety}; the notions of inclusive variety and its generation are not defined in the supplied record.

\textbf{Classification problem (G. Mashevitzky).} Classify semigroups from the basic classes of semigroups up to inclusive equivalence. Once an intended basic class \(\mathcal K\) and the meaning of identical inclusion have been specified, the task is to describe the equivalence classes
\[
[S]_{\mathcal K,\mathrm{inc}}
=\{T\in\mathcal K:\operatorname{Inc}(T)=\operatorname{Inc}(S)\},
\qquad S\in\mathcal K.
\]
The supplied record does not identify the intended basic classes, specify whether comparisons across different basic classes are required, or state a convention concerning empty semigroups. These qualifications remain unresolved; the formulation above records the requested classification relative to the missing source definitions.
\subsection{Short English statement}
Classify semigroups in the source's intended basic classes according to which identical inclusions they satisfy. The supplied record attributes the problem to G. Mashevitzky but does not define identical inclusions, inclusive varieties, or the intended basic classes, so a fully self-contained classification target requires further source clarification.
\subsection{Sources}
\begin{itemize}
\item Local UnsolvedMath registry: permanent ID \(800040\), exact source problem code \texttt{catalogue entry}, title \emph{The Lyapin's Notebook --- Problem 42}; supplied upstream revision \texttt{d371ee917ac9647fbb06cd6fc529deec6f3dda82}.
\item Supplied source reference: Bershadsky, Kublanovsky, Mashevitzky, \emph{The Lyapin's Notebook: Unsolved Problems in Semigroup Theory} (2026), Problem 42, \url{https://arxiv.org/abs/2604.04763}. The record reports access on 2026-10-16 and extraction by \texttt{arxiv\_latex}; these are supplied provenance details, not an independent source check.
\item The statement explicitly attributes the problem to G. Mashevitzky. The registry separately lists Bershadsky, Kublanovsky, Mashevitzky in its proposer field and gives 2026 as the proposed year; the supplied evidence does not resolve whether that year denotes original proposal or collection publication.
\item The registry labels the problem \emph{open}. This is a supplied status label, not an independently verified mathematical or literature assessment.
\end{itemize}
\end{document}
